\section{Introduction}

Bitcoin has evolved from a largely retail and exchange-specific trading venue into a globally traded asset held, intermediated, and hedged through increasingly institutional channels. The approval and subsequent trading of spot Bitcoin exchange-traded products have made Bitcoin exposure accessible through conventional portfolio infrastructure, while the development of liquid listed and over-the-counter option markets has created a cross-section of prices that embeds forward-looking assessments of volatility, jump risk, and tail-state compensation. Yet the economic object implied by these option prices remains only partially understood. Bitcoin is traded continuously across venues and jurisdictions, and its return dynamics display precisely the features that make standard asset-pricing restrictions difficult to impose: strong non-Gaussianity, heavy tails, abrupt price discontinuities, volatility clustering, and recurrent regime shifts. In such an environment, neither a representative-agent Euler equation estimated from low-frequency consumption data nor a lognormal diffusion calibrated to a single volatility parameter is likely to provide a satisfactory account of how investors value Bitcoin payoffs.

{\color{red}This paper studies Bitcoin option prices through the lens of the pricing kernel. For a payoff $\psi(S_T)$ written on the terminal Bitcoin price $S_T$, no-arbitrage valuation gives
\begin{equation}
P_t=E_t^\mathbb{Q}\left[e^{-r\tau}\psi(S_T)\right]
    =E_t^\mathbb{P}\left[e^{-r\tau}K_t(S_T)\psi(S_T)\right],
\end{equation}
where $\mathbb{P}$ is the physical measure, $\mathbb{Q}$ is the risk-neutral measure, $\tau=T-t$, and
\begin{equation}    
K_t(x)=\frac{q_t(x)}{p_t(x)}
\end{equation}
is the ratio of the risk-neutral density to the physical density.} This object is central because it converts statistical probabilities into state prices. A dollar paid in a state with high marginal value should command a larger state price than a dollar paid in a state with low marginal value; consequently, the shape of $K_t$ summarizes how the market prices downside risk, upside speculation, variance risk, and jump risk. The pricing-kernel representation also makes clear why Bitcoin options are an informative laboratory: the same terminal distribution may receive very different physical and risk-neutral weights when investors demand compensation for crash risk or for volatility-jump exposure.

The classical stochastic-discount-factor literature provides the first benchmark for this problem. \citet{HansenSingleton1982} estimate Euler-equation restrictions by generalized method of moments, showing how asset returns can discipline intertemporal marginal rates of substitution. Hansen and Jagannathan \citep{HansenJagannathan1991} derive volatility bounds for admissible stochastic discount factors, thereby turning asset-return data into a diagnostic restriction on pricing models. These approaches are foundational, but they are not directly tailored to Bitcoin. Euler-equation estimation requires specifying preferences and information variables, often relying on consumption or macroeconomic data that are difficult to interpret for a globally traded cryptoasset. Hansen--Jagannathan bounds are powerful diagnostics, but they do not by themselves recover the state-by-state pricing kernel embedded in option prices. Our setting therefore calls for an option-based approach that can recover the relative valuation of future Bitcoin states more directly.

A parallel option-based literature exploits the fact that option prices reveal state prices. Breeden and Litzenberger \citep{BreedenLitzenberger1978} show that the second derivative of call prices with respect to strike identifies state-contingent prices. Building on this insight, Shimko \citep{Shimko1993} smooths implied volatility before recovering the risk-neutral density, while A{\"i}t-Sahalia and Lo \citep{AitSahaliaLo1998} develop a nonparametric estimator of state-price densities implicit in option prices. A{\"i}t-Sahalia and Lo \citep{AitSahaliaLo2000} further combine option-implied state prices with estimates of the physical distribution to infer implied risk aversion. These studies formalize the empirical density-ratio logic underlying this paper: once the physical density $p_t$ and risk-neutral density $q_t$ are estimated for the same horizon and state variable, their ratio gives an empirical pricing kernel. In practice, however, both densities are difficult to estimate. Risk-neutral density recovery is a continuous-strike identification result applied to discrete, noisy, maturity-specific option data; physical density estimation is sensitive to the historical window, tail observations, and the assumed return dynamics. These difficulties are especially acute for Bitcoin, where tail states are economically central rather than peripheral.

The empirical pricing-kernel literature also warns that naively estimated density ratios can produce economically puzzling shapes. Jackwerth \citep{Jackwerth2000} and A{\"i}t-Sahalia and Lo \citep{AitSahaliaLo2000} document pricing kernels that are locally increasing over parts of the state space, a finding commonly referred to as the pricing-kernel puzzle because it appears inconsistent with a globally decreasing marginal utility schedule. Rosenberg and Engle \citep{RosenbergEngle2002} estimate empirical pricing kernels using flexible physical densities and parametric restrictions on the kernel, emphasizing the role of the physical return model in shaping the density ratio. Christoffersen, Heston, and Jacobs \citep{ChristoffersenHestonJacobs2013} show that option anomalies can be captured by a variance-dependent pricing kernel, while Linn, Shive, and Shumway \citep{LinnShiveShumway2018} argue that conditioning and information-set mismatch are central to understanding option-implied stochastic discount factors. These contributions imply that a Bitcoin pricing kernel should not be estimated as an unconditional object averaged across market regimes. It should be conditional on the current volatility state, option-implied distribution, and market environment.

The Bitcoin setting sharpens these econometric and economic issues. First, the physical density is difficult to estimate. Historical kernel density estimators are flexible but fragile in the tails; GARCH-type models capture volatility clustering but typically handle jumps only indirectly; and purely diffusion-based stochastic-volatility models can understate discontinuities in returns and volatility. Second, the risk-neutral density is also difficult to recover. Breeden--Litzenberger identification requires a smooth continuum of strikes, whereas observed Bitcoin options are discrete, liquidity-constrained, and unevenly distributed across moneyness and maturities. Nonparametric smoothing of the implied-volatility surface reduces noise but introduces bandwidth, interpolation, extrapolation, and no-arbitrage choices. Third, the density-ratio construction magnifies errors: small mistakes in either $p_t$ or $q_t$ can become economically large when $q_t/p_t$ is evaluated in the tails. These concerns motivate a model that is sufficiently structured to stabilize tail inference, but sufficiently rich to match the discontinuous and heteroskedastic nature of Bitcoin returns.

We therefore estimate Bitcoin pricing kernels under the stochastic volatility with correlated jumps (SVCJ) model. The SVCJ framework extends stochastic-volatility models by allowing both the return and variance processes to jump, with contemporaneous dependence between return-jump and volatility-jump sizes. This feature is economically important for Bitcoin: large negative returns often coincide with abrupt increases in uncertainty, and option prices should reflect compensation not only for diffusive variance risk but also for jump-arrival risk and co-jump risk. Our model builds on the jump-diffusion and stochastic-volatility option-pricing literature, including Merton \citep{Merton1976}, Heston \citep{Heston1993}, Duffie, Pan, and Singleton \citep{DuffiePanSingleton2000}, and Eraker, Johannes, and Polson \citep{ErakerJohannesPolson2003}. Recent cryptocurrency-option studies similarly emphasize that Bitcoin option prices require models with stochastic volatility, jumps, and co-jumps \citep{HouWangChenHaerdle2020,MaticPackhamHaerdle2021}. Relative to these studies, our focus is not only on option valuation, but also on the state-price implications of the model through the pricing kernel.

Our empirical design uses historical Bitcoin returns to estimate the physical distribution and Deribit Bitcoin options to estimate the risk-neutral distribution, imposing the same SVCJ structure under both measures. This common dynamic structure is important. It allows differences between $P$ and $Q$ to be interpreted as compensation for variance risk, jump risk, and return-volatility co-jump risk rather than as artifacts of using unrelated density estimators. It also allows the pricing kernel to be conditional: for each estimation date and maturity, the kernel is tied to the prevailing level of volatility, the calibrated option surface, and the current state of the Bitcoin market. In contrast to an unconditional density ratio that averages across market regimes, our approach is designed to ask whether Bitcoin state prices behave differently during calm, speculative, and stress periods.

This paper makes three contributions. First, it contributes to the empirical pricing-kernel literature by constructing a robust, conditional pricing kernel for a nontraditional asset whose trading calendar, investor base, and tail behavior differ sharply from equity-index markets. Second, it contributes to cryptocurrency option pricing by using a unified SVCJ framework to compare the physical and risk-neutral laws of Bitcoin, thereby quantifying variance, jump, and co-jump risk premia embedded in option prices. Third, it contributes to the literature on the pricing-kernel puzzle by examining whether the nonmonotonicities documented in index-option markets also arise in Bitcoin, and whether they are robust to a model that explicitly accounts for stochastic volatility and correlated jumps. Recent evidence suggests that Bitcoin risk premia are state-dependent and that Bitcoin pricing kernels can exhibit pronounced nonlinear shapes \citep{AlmeidaGrithMiftachovWang2024}; our framework provides a structural way to study these patterns through the separate channels of diffusive volatility, jump intensity, and jump-size dependence.

The rest of the paper proceeds as follows. Section~2 introduces data. Section~3 introduces the SVCJ model under the physical and risk-neutral measures. Section~4 describes the Bitcoin return and option data and the estimation strategy. Section~5 presents the estimated physical and risk-neutral distributions, the resulting pricing kernels, and robustness checks across maturities and market regimes. Section~6 discusses implications for Bitcoin risk premia, option valuation, and portfolio allocation. Section~7 concludes.

